% Generated draft: requires mathematical and source review.
\documentclass[11pt]{article}
\usepackage[margin=25mm]{geometry}
\usepackage{fontspec}
\setmainfont{Latin Modern Roman}
\usepackage{amsmath,amssymb,mathtools,unicode-math}
\setmathfont{Latin Modern Math}
\usepackage{enumitem,needspace,xurl,hyperref}
\setlength{\parindent}{0pt}
\setlength{\parskip}{5pt}
\setlength{\emergencystretch}{3em}
\newcommand{\N}{\mathbb{N}}
\newcommand{\Z}{\mathbb{Z}}
\newcommand{\Q}{\mathbb{Q}}
\newcommand{\R}{\mathbb{R}}
\newcommand{\C}{\mathbb{C}}
\newcommand{\F}{\mathbb{F}}
\newcommand{\A}{\mathbb{A}}
\newcommand{\PP}{\mathbb{P}}
\newcommand{\E}{\mathbb{E}}
\newcommand{\eps}{\varepsilon}
\newcommand{\dd}{\,\mathrm{d}}
\begin{document}
\section*{UnsolvedMath ID: 30000394}
\textbf{Source problem code:} catalogue entry
\subsection{Definitions and mathematical statement}
\textbf{Ore-type packing conjecture, as supplied.} Let \(G_1,G_2\) be finite undirected simple graphs on \(n\ge1\) vertices. For a graph with at least one edge define\[\sigma_2(G)=\max_{uv\in E(G)}\frac{\deg_G(u)+\deg_G(v)}2,\qquad \Delta(G)=\max_{v\in V(G)}\deg_G(v).\] If\[(\sigma_2(G_1)+1)(\Delta(G_2)+1)\le n+1,\] must \(G_1\) and \(G_2\) pack? Packing means that there is a bijection \(\varphi:V(G_1)\to V(G_2)\) for which \(uv\in E(G_1)\) implies \(\varphi(u)\varphi(v)\notin E(G_2)\). Equivalently, both graphs can be placed on the vertex set of \(K_n\) with disjoint edge sets.\par The source's \(\sigma_2\) is exactly a maximum over adjacent vertex pairs divided by two. It is not a minimum over nonadjacent pairs, despite other uses of Ore-type notation, and half-integral values are not rounded. The hypothesis uses \(\sigma_2\) for \(G_1\) and maximum degree for \(G_2\); it is not replaced by a symmetric product of two \(\sigma_2\) parameters. If \(G_1\) is edgeless, packing is automatic but the displayed maximum is undefined. One may handle that case separately, as done here, rather than silently assigning a value to the empty maximum. The source does not explicitly state graph simplicity or an empty-edge convention.\par The supplied triage checked 2026-08-21 labels the item open, calls it an Ore-type BEC analogue, and says the report presents it as a conjecture without resolving it. This describes the supplied report and is not independent verification of its current literature status.
\subsection{Short English statement}
The conjecture gives a sufficient product bound for relabeling two graphs so that their edges do not overlap. Its degree parameter is the maximum average degree sum across an edge, exactly as supplied. The edgeless case is separated because the defining maximum is otherwise empty.
\subsection{Sources}
\begin{itemize}\item Local UnsolvedMath registry, permanent ID 30000394, code catalogue entry; \emph{Ore-Type Conditions for Graph Packing}, Oberwolfach Reports Open Problems.\item \emph{Combinatorics (2006)}, Oberwolfach Reports, DOI 10.4171/owr/2006/01, pp. 5--72. \url{https://doi.org/10.4171/owr/2006/01}. Proposed year 2006; no proposer named.\item Supplied literature label \emph{original Oberwolfach report}: \url{https://doi.org/10.4171/owr/2006/01}. No separate authoritative conjecture citation is supplied.\item Literature triage checked 2026-08-21, status \texttt{open}; record updated 2026-09-16. The \texttt{corrected\_verified} field reports source checking, not independently verified here.\end{itemize}
\end{document}
